\subsection{Excluding finite mass-shell profiles}
\label{sec:trace-profiles}
The centered trace has a short-distance behavior that can distinguish it from
finite polynomial expressions in free fields. The relevant distinction is the
negative power of the logarithm. We prove the exclusion directly for the
mass-shell profiles appearing in the statement. The proof proceeds through
finite power--logarithm approximations to Laplace integrals; it does not
assume a convergent asymptotic series.

\subsubsection{An algebra of finite asymptotic approximations}
Let $\mathcal P$ be the real vector space of finite sums of the functions
$s^a(\log s)^k$, where $a\in\mathbb Z$, $k\in\mathbb N$ and $s>0$.
It is an algebra containing the constants. Say that $f$ has an approximation
of order $b\in\mathbb Z$ if
\[
 f(s)=P(s)+O(s^b)\quad(s\downarrow0),\qquad P\in\mathcal P.
\]
Say that $f$ has power--logarithm approximations if it has an approximation
of every integer order. Different orders may use different finite sums.
An approximation of order $b$ is also one of any order $a\leq b$, since
$s^b\leq s^a$ for $0<s<1$.

\begin{lemma}\label{lem:profile-algebra}
Functions with power--logarithm approximations form a real algebra. Every
such function, and every element of $\mathcal P$, is $O(s^u)$ for some
integer $u$. Multiplication by $s^a$ shifts the order of an approximation
by $a$.
\end{lemma}
\begin{proof}
For $0<s<1$, $|\log s|\leq s^{-1}$. Thus each monomial is bounded by
$s^{a-k}$, and a finite sum is $O(s^u)$ for the smallest of the finitely
many exponents $a-k$. An order-zero approximation then gives the growth
claim for $f$, after replacing $u$ by $\min(u,0)$. The assertions about
addition, real scalar multiplication, and shifts are immediate from the
definition and closure of $\mathcal P$ under these operations.

For multiplication fix a desired order $b$. Choose $\ell$ with
$g=O(s^\ell)$ and an approximation $f=P+O(s^{b-\ell})$. Choose $u$ with
$P=O(s^u)$ and then $g=Q+O(s^{b-u})$. The exact identity
\[
 fg-PQ=(f-P)g+P(g-Q)
\]
shows that both terms are $O(s^b)$. Since $PQ\in\mathcal P$, this proves
the required approximation. In particular it proves closure under every
finite product, including the empty product $1$.
\end{proof}

\subsubsection{Radial Laplace integrals}
For $m\geq0$ and $j,k\in\mathbb N$, put
\[
 M_{j,k}(s)=\int_0^\infty r^{j+2}(m^2+r^2)^{(k-1)/2}
                         e^{-s\sqrt{m^2+r^2}}\,dr,\qquad s>0.
\]
\begin{proposition}\label{prop:profile-moments}
Every $M_{j,k}$ has power--logarithm approximations.
\end{proposition}
\begin{proof}
If $m=0$ the integral is $(j+k+1)!s^{-(j+k+2)}$. If $m>0$, set
$t=\sqrt{1+(r/m)^2}$ and $\nu=(j+2)/2$. The integral for $k=0$
becomes $m^{j+2}\int_1^\infty e^{-mst}(t^2-1)^{\nu-1/2}\,dt$.
Since $ms>0$ and $\nu>-1/2$, the Bessel integral
\cite[Eq.~10.32.8]{DLMF} gives
\[
 M_{j,0}(s)=\frac{m^{j+2}\Gamma(\nu+1/2)}{\sqrt\pi}
                    \left(\frac{2}{ms}\right)^\nu K_\nu(ms).
\]
For integer $\nu$, Eq.~10.31.1 of \cite{DLMF} writes
$z^{-\nu}K_\nu(z)$ as a Laurent series with finitely many negative
powers plus $\log z$ times a power series. For noninteger $\nu$,
Eqs.~10.25.2 and 10.27.4 express it as a power series plus
$z^{-2\nu}$ times a power series. Here $2\nu=j+2$ is integral.
To control derivatives of any fixed order, truncate the analytic
factors sufficiently far. Differentiating an analytic tail $s^a h(s)$
lowers its order by at most the number of derivatives. A log factor
costs at most one further power in the bound.
Replacing $z$ by $ms$ changes $\log z$ to $\log s+\log m$.
Finally $M_{j,k}=(-\partial_s)^k M_{j,0}$ by dominated differentiation
on compact subsets of $s>0$, since $m>0$ and the differentiated
integrands decay exponentially at infinity. This supplies the
approximations of every integer order.
\end{proof}

\subsubsection{From modes to the logarithmic obstruction}
A mode consists of a mass $m\geq0$, integers $j,k\geq0$, and a continuous
real function $P$ on $\mathbb R^3$ homogeneous of degree $j$ under positive
scalings. Its Euclidean profile is
\begin{equation}\label{eq:profile-mode}
 K_{m,P,k}(s)=\int_{\mathbb R^3}
       \frac{\omega_m(p)^kP(p)}{2\omega_m(p)}e^{-s\omega_m(p)}\,dp,
 \qquad \omega_m(p)=\sqrt{m^2+|p|^2}.
\end{equation}
Let $d\sigma$ denote ordinary surface area measure on the unit sphere.
The maximum of $|P|$ on that sphere is finite. In polar coordinates,
$|P(r\theta)|\leq Cr^j$, so absolute integrability follows from that
of $M_{j,k}$. Fubini's theorem and homogeneity therefore give the exact
normalization
\begin{equation}\label{eq:profile-angular}
 K_{m,P,k}(s)=\frac12\left(\int_{S^2}P(\theta)\,d\sigma(\theta)\right)
                         M_{j,k}(s).
\end{equation}
A track is a finite product of mode profiles; the empty track is $1$.
By Proposition~\ref{prop:profile-moments} and
Lemma~\ref{lem:profile-algebra}, every track has power--logarithm
approximations. The same is true of the real part of a finite complex
linear combination of tracks, because all the individual tracks are real.

\begin{lemma}[No inverse logarithmic leading term]\label{lem:profile-log}
If a real function $f$ has an approximation of order zero and
\[
 \lim_{s\downarrow0}s^8\bigl(\log s^{-1}\bigr)^2f(s)=A\in\mathbb R,
\]
then $A=0$.
\end{lemma}
\begin{proof}
Write $f=P+O(1)$ with $P\in\mathcal P$. Since
$s^8(\log s^{-1})^2\to0$, the same weighted limit holds for $P$.
Restrict to $s_n=2^{-n}$. A monomial $c_{a,k}s^a(\log s)^k$ contributes
\[
 c_{a,k}\bigl(\log(1/2)\bigr)^{k+2}
             \bigl(2^{-(8+a)}\bigr)^n n^{k+2}.
\]
Combine equal pairs $(a,k)$ and discard zero coefficients. The remaining
finite sum has distinct pairs $(\lambda,q)=(2^{-(8+a)},k+2)$,
with $\lambda>0$ and $q\geq2$. Order them lexicographically and choose
the largest pair $(\lambda_*,q_*)$. If the sum is empty, its limit is
zero. If $\lambda_*<1$, every term tends to zero: writing
$\lambda=e^{-c}$ with $c>0$ reduces this to $n^qe^{-cn}\to0$.
For completeness, $e^{cn}\geq(cn)^{q+1}/(q+1)!$ proves that limit.

If $\lambda_*\geq1$, then $D_n=\lambda_*^n n^{q_*}\to\infty$.
Each smaller term divided by $D_n$ tends to zero. With a smaller base,
exponential decay dominates any remaining power of $n$ by the same
estimate; with equal base, the ratio is a negative power of $n$.
Thus the finite sum divided by $D_n$ tends to its nonzero leading
coefficient. But the sum has the finite limit $A$, and division by
$D_n\to\infty$ makes its limit zero, a contradiction. Only the first
case is possible, and it has $A=0$.
\end{proof}

\begin{theorem}\label{thm:profile-exclusion}
Let $T:(0,\infty)\to\mathbb C$ satisfy
\[
 \lim_{s\downarrow0}s^8\bigl(\log s^{-1}\bigr)^2T(s)=A,
 \qquad A>0.
\]
Then $T$ cannot agree on any interval $(0,\delta)$ with a finite complex
linear combination of tracks built from \eqref{eq:profile-mode}.
\end{theorem}
\begin{proof}
If such agreement held, the real part of that combination would have an
order-zero approximation. Taking real parts in the displayed limit and
using the agreement would contradict Lemma~\ref{lem:profile-log}.
\end{proof}
This proves the exclusion once the positive trace asymptotic has been
established for the constructed theory. It imposes no restriction on an
infinite sum or a continuous distribution of masses, and no such exclusion
is asserted here. The construction of $T$ and its positive coefficient $A$
is a separate quantitative step.
